\begin{ex}
Một hình chữ nhật có chiều dài 8 cm, chiều rộng 5 cm. Tính diện tích hình chữ nhật đó theo xăng-ti-mét vuông.
\shortans{40}
\loigiai{
$8 \times 5 = 40\text{ cm}^2$.
}
\end{ex}

\begin{ex}
Một hình chữ nhật có chiều dài 9 cm, chiều rộng 6 cm. Tính diện tích hình chữ nhật đó theo xăng-ti-mét vuông.
\shortans{54}
\loigiai{
$9 \times 6 = 54\text{ cm}^2$.
}
\end{ex}

\begin{ex}
Một hình chữ nhật có chiều dài 12 cm, chiều rộng 4 cm. Diện tích của hình chữ nhật đó là bao nhiêu xăng-ti-mét vuông?
\shortans{48}
\loigiai{
$12 \times 4 = 48\text{ cm}^2$.
}
\end{ex}

\begin{ex}
Một hình vuông có cạnh dài 6 cm. Tính diện tích hình vuông đó theo xăng-ti-mét vuông.
\shortans{36}
\loigiai{
$6 \times 6 = 36\text{ cm}^2$.
}
\end{ex}

\begin{ex}
Một hình vuông có cạnh dài 8 cm. Tính diện tích hình vuông đó theo xăng-ti-mét vuông.
\shortans{64}
\loigiai{
$8 \times 8 = 64\text{ cm}^2$.
}
\end{ex}

\begin{ex}
Một hình vuông có cạnh dài 9 cm. Tính diện tích hình vuông đó theo xăng-ti-mét vuông.
\shortans{81}
\loigiai{
$9 \times 9 = 81\text{ cm}^2$.
}
\end{ex}

\begin{ex}
Một hình chữ nhật có diện tích $35\text{ cm}^2$, chiều dài là 7 cm. Chiều rộng của hình chữ nhật đó là bao nhiêu xăng-ti-mét?
\shortans{5}
\loigiai{
$35 : 7 = 5$ cm.
}
\end{ex}

\begin{ex}
Một hình chữ nhật có diện tích $48\text{ cm}^2$, chiều rộng là 6 cm. Chiều dài của hình chữ nhật đó là bao nhiêu xăng-ti-mét?
\shortans{8}
\loigiai{
$48 : 6 = 8$ cm.
}
\end{ex}

\begin{ex}
Một hình vuông có chu vi 24 cm. Cạnh của hình vuông đó dài bao nhiêu xăng-ti-mét?
\shortans{6}
\loigiai{
$24 : 4 = 6$ cm.
}
\end{ex}

\begin{ex}
Một hình vuông có chu vi 28 cm. Tính diện tích hình vuông đó theo xăng-ti-mét vuông.
\shortans{49}
\loigiai{
Cạnh = $28 : 4 = 7$ cm. Diện tích = $7 \times 7 = 49\text{ cm}^2$.
}
\end{ex}

\begin{ex}
Một hình vuông có chu vi 36 cm. Diện tích hình vuông đó là bao nhiêu xăng-ti-mét vuông?
\shortans{81}
\loigiai{
Cạnh = $36 : 4 = 9$ cm. Diện tích = $9 \times 9 = 81\text{ cm}^2$.
}
\end{ex}

\begin{ex}
Một hình vuông có chu vi 20 cm. Diện tích hình vuông đó là bao nhiêu xăng-ti-mét vuông?
\shortans{25}
\loigiai{
Cạnh = $20 : 4 = 5$ cm. Diện tích = $5 \times 5 = 25\text{ cm}^2$.
}
\end{ex}

\begin{ex}
Một hình chữ nhật có chiều rộng 4 cm, chiều dài gấp 3 lần chiều rộng. Diện tích hình chữ nhật đó là bao nhiêu xăng-ti-mét vuông?
\shortans{48}
\loigiai{
Chiều dài = $4 \times 3 = 12$ cm. Diện tích = $12 \times 4 = 48\text{ cm}^2$.
}
\end{ex}

\begin{ex}
Một hình chữ nhật có nửa chu vi là 15 cm, chiều dài là 9 cm. Diện tích hình chữ nhật đó là bao nhiêu xăng-ti-mét vuông?
\shortans{54}
\loigiai{
Chiều rộng = $15 - 9 = 6$ cm. Diện tích = $9 \times 6 = 54\text{ cm}^2$.
}
\end{ex}

\begin{ex}
Một tấm thiệp hình chữ nhật có chiều dài 15 cm, chiều rộng 10 cm. Diện tích tấm thiệp là bao nhiêu xăng-ti-mét vuông?
\shortans{150}
\loigiai{
$15 \times 10 = 150\text{ cm}^2$.
}
\end{ex}

\begin{ex}
Một viên gạch men hình vuông có cạnh dài 20 cm. Diện tích viên gạch men là bao nhiêu xăng-ti-mét vuông?
\shortans{400}
\loigiai{
$20 \times 20 = 400\text{ cm}^2$.
}
\end{ex}

\begin{ex}
Một hình vuông có diện tích $16\text{ cm}^2$. Độ dài cạnh của hình vuông đó là bao nhiêu xăng-ti-mét?
\shortans{4}
\loigiai{
Vì $4 \times 4 = 16$ nên cạnh bằng 4 cm.
}
\end{ex}

\begin{ex}
Một hình vuông có diện tích $100\text{ cm}^2$. Cạnh hình vuông dài bao nhiêu xăng-ti-mét?
\shortans{10}
\loigiai{
Vì $10 \times 10 = 100$ nên cạnh bằng 10 cm.
}
\end{ex}

\begin{ex}
Một hình chữ nhật có chu vi 30 cm, chiều dài 10 cm. Diện tích hình chữ nhật đó là bao nhiêu xăng-ti-mét vuông?
\shortans{50}
\loigiai{
Nửa chu vi = 15 cm; rộng = $15 - 10 = 5$ cm. Diện tích = $10 \times 5 = 50\text{ cm}^2$.
}
\end{ex}

\begin{ex}
Một hình chữ nhật có chiều dài 14 cm, chiều rộng bằng nửa chiều dài. Diện tích hình chữ nhật đó là bao nhiêu xăng-ti-mét vuông?
\shortans{98}
\loigiai{
Chiều rộng = 7 cm. Diện tích = $14 \times 7 = 98\text{ cm}^2$.
}
\end{ex}

\begin{ex}
Ghép hai hình vuông cạnh 4 cm thành một hình chữ nhật. Diện tích hình chữ nhật tạo thành là bao nhiêu xăng-ti-mét vuông?
\shortans{32}
\loigiai{
Diện tích 1 hình vuông = $4 \times 4 = 16\text{ cm}^2$. Hai hình = $16 \times 2 = 32\text{ cm}^2$.
}
\end{ex}

\begin{ex}
Một miếng bìa hình chữ nhật dài 18 cm, rộng 5 cm. Người ta cắt ra một hình vuông cạnh 5 cm. Diện tích phần bìa còn lại là bao nhiêu xăng-ti-mét vuông?
\shortans{65}
\loigiai{
Diện tích ban đầu = $18 \times 5 = 90\text{ cm}^2$. Cắt đi = $5 \times 5 = 25\text{ cm}^2$. Còn lại: $90 - 25 = 65\text{ cm}^2$.
}
\end{ex}

\begin{ex}
Nếu gấp cạnh của hình vuông lên 2 lần thì diện tích hình vuông gấp lên mấy lần?
\shortans{4}
\loigiai{
Diện tích tăng lên $2 \times 2 = 4$ lần.
}
\end{ex}

\begin{ex}
Một khung ảnh hình vuông cạnh 12 cm. Diện tích mặt kính của khung ảnh là bao nhiêu xăng-ti-mét vuông?
\shortans{144}
\loigiai{
$12 \times 12 = 144\text{ cm}^2$.
}
\end{ex}

\begin{ex}
Tính diện tích hình chữ nhật có chiều dài 20 cm và chiều rộng 7 cm theo xăng-ti-mét vuông.
\shortans{140}
\loigiai{
$20 \times 7 = 140\text{ cm}^2$.
}
\end{ex}
